\centerline{\bf \Huge Первый разнобой 2025}

\begin{flushright}
        \scriptsize{ Если мои ответы пугают тебя, перестань задавать страшные вопросы.\\Джулс Уиннфилд. Криминальное чтиво}
\end{flushright}

\problem{1}
Андрей взял деревянный кубик и написал на каждой грани натуральное число от 1 до 6. Денис заметил, что каждое из этих чисел делит сумму чисел, стоящих на соседних по стороне гранях. Могло ли случиться так, что все шесть чисел различны?

\problem{2}
Положительные числа $x$ и $у$ таковы, что $(1+x)(1+y)=2$. Докажите, что

$$
x y+\frac{1}{x y} \geqslant 6
$$

\problem{3}
В Первой стране есть 32 города, и во Второй стране тоже 32 города. Между некоторыми городами налажено прямое двустороннее авиасообщение. Известно, что из каждого города Первой страны можно напрямую долететь хотя бы до 26 городов Второй страны. Докажите, что найдутся два города Второй страны таких, что в любой город Первой страны можно напрямую долететь хотя бы из одного из них.

\problem{4}
Диагонали трапеции $A B C D\ (A D \| B C)$ пересекаются в точке $O$. На стороне $A D$ выбрали точку $X$. Из вершины $A$ опустили перпендикуляр $A S$ на прямую $C X$, а из вершины $D$ - перпендикуляр $D T$ на прямую $B X$. Оказалось, что длины этих перпендикуляров равны. Докажите, что $\angle O X B=\angle O X C$.

\problem{5}
Компьютер вычисляет суммы последовательных простых чисел, начиная с 2 . Результаты вычислений выглядят следующим образом:

--- сначала вычисляется $2+3$;

--- затем $-2+3+5$;

--- затем $-2+3+5+7$;

--- и так далее.

Могут ли две подряд идущие вычисленные суммы одновременно быть точными квадратами?

